\chapter{Background}\label{cap:background}
\glsresetall

% =====================================================================
% SKELETON. Target: 12--18 pages.
%
% Purpose: give the banca exactly the machinery needed to read Cap04, and
% nothing more. Every section here must be load-bearing for the method or
% the evaluation. If a section does not get referenced later, cut it.
%
% Verb tense for literature review: simple present ("X proposes Y").
% =====================================================================

This chapter introduces the machinery needed to read the method and evaluation that follow, and
no more. It proceeds from the generative backbone outward: diffusion models
(Section~\ref{sec:bg-diffusion}) and how they are conditioned (Section~\ref{sec:bg-conditioning});
the temporal mechanism that turns a still-image model into a clip generator
(Section~\ref{sec:bg-temporal}); and the camera-parameter conditioning of Generative Photography
(Section~\ref{sec:bg-genphoto}), whose architecture this work reuses. It then turns to the
expression side: the Facial Action Coding System (Section~\ref{sec:bg-facs}), the tools that
estimate action units (Section~\ref{sec:bg-au-estimation}), the datasets that supply
intensity-graded expressions (Section~\ref{sec:bg-datasets}), and the metrics used to evaluate
control (Section~\ref{sec:bg-metrics}).

% =====================================================================
\section{Diffusion Models}\label{sec:bg-diffusion}

% Forward/reverse process, the noise-prediction objective, and why latent
% space matters for tractability. Keep the derivation minimal: state the
% training objective and move on.
%   \citep{ho2020ddpm}      -- DDPM, forward/reverse formulation
%   \citep{song2021ddim}    -- deterministic sampling, fewer steps
%   \citep{rombach2022ldm}  -- latent diffusion / Stable Diffusion, the backbone here

A diffusion model \citep{ho2020ddpm} defines a forward process that gradually corrupts a data
sample $x_0$ into noise over $T$ steps, $x_t = \sqrt{\bar\alpha_t}\,x_0 +
\sqrt{1-\bar\alpha_t}\,\epsilon$ with $\epsilon\sim\mathcal{N}(0,I)$, and learns to reverse it. In
practice the network $\epsilon_\theta$ is trained to predict the added noise, minimising
$\mathbb{E}_{x_0,\epsilon,t}\,\lVert \epsilon - \epsilon_\theta(x_t,t) \rVert^2$; generation then
denoises from pure noise, and DDIM \citep{song2021ddim} makes this deterministic and possible in
few steps. Operating in pixel space is costly, so latent diffusion \citep{rombach2022ldm} runs the
process in the compressed latent space of a pretrained \gls{vae}, whose encoder maps images to
latents for training and whose decoder maps generated latents back to pixels. Stable Diffusion,
the backbone used in this work, is a latent diffusion model of exactly this form, and every
conditioning mechanism below acts on its denoising network.

% =====================================================================
\section{Conditioning Mechanisms in Text-to-Image Diffusion}\label{sec:bg-conditioning}

% This is the section Cap04 depends on most. Cover, in this order:
%   1. Text conditioning via cross-attention with a CLIP text encoder
%      \citep{radford2021clip}; why it is categorical rather than continuous.
%   2. Classifier-free guidance \citep{ho2022cfg} and its use on a
%      non-text condition.
%   3. Adapter-style injection on a FROZEN backbone: ControlNet
%      \citep{zhang2023controlnet}, IP-Adapter \citep{ye2023ipadapter},
%      LoRA \citep{hu2022lora}. Establish the vocabulary of "decoupled
%      cross-attention" and "zero-initialised merge", both used in Cap04.

Text-to-image diffusion conditions the denoising network on a prompt encoded by a CLIP text
encoder \citep{radford2021clip} and injected through cross-attention layers, in which the image
features attend to the text tokens. This conditioning is effective but \emph{categorical}: a prompt
can say ``smiling'' or ``not smiling'', but expressing \emph{how much} of an expression is present,
on a continuous scale, is beyond what a text token cleanly encodes. Classifier-free guidance
\citep{ho2022cfg} sharpens adherence to any condition by extrapolating between the conditional and
unconditional predictions at sampling time, and applies to non-text conditions as well.

Rather than fine-tuning the whole backbone, control is increasingly added by small modules on a
\emph{frozen} network. ControlNet \citep{zhang2023controlnet} attaches a trainable copy of the
encoder for spatial conditions; IP-Adapter \citep{ye2023ipadapter} adds an image condition through
\emph{decoupled cross-attention}, a parallel attention branch whose output is summed into the text
branch; and \gls{lora} \citep{hu2022lora} injects low-rank updates into existing weight matrices.
Two ideas from this literature recur in Chapter~\ref{cap:method}: decoupled cross-attention, and
the \emph{zero-initialised merge} -- initialising an injected branch to zero so that, before
training, the modified model is identical to the frozen original.

% =====================================================================
\section{Temporal Generation and Motion Modules}\label{sec:bg-temporal}

% The mechanism that makes this work's claim possible: a motion module
% inserted into a 2D UNet turns it into a 3D UNet whose temporal attention
% binds the frames of one clip to a shared identity and scene.
%   \citep{guo2024animatediff} -- AnimateDiff motion module
% Make explicit the point the whole contribution rests on: with temporal
% attention, the frames of a clip are generated JOINTLY, whereas N
% independent stills share only a random seed.

A still-image diffusion model generates each image independently. AnimateDiff
\citep{guo2024animatediff} makes it generate short clips by inserting a \emph{motion module} into
the 2D UNet: temporal-attention layers that let the feature maps of different frames attend to one
another, turning the 2D network into a 3D one. The consequence is central to this work. Because the
frames of a clip are produced \emph{jointly} through temporal attention, they share a single
identity and scene by construction. This is exactly what distinguishes a generated clip from $N$
images sampled independently under a shared random seed: the latter are correlated only through
their starting noise, whereas temporal attention binds them architecturally. The prior art of
Chapter~\ref{cap:related-work} produces the former; this work produces the latter.

% =====================================================================
\section{Parametric Control of Generative Cameras}\label{sec:bg-genphoto}

% The direct antecedent, presented here as background rather than as
% related work because the method in Cap04 reuses its mechanism verbatim.
%   \citep{yuan2024genphoto}
% Cover: the 6-channel physical conditioning layout, the contrastive camera
% learning (\gls{ccl}) text encoding that represents *change between frames*,
% and the zero-initialised merge into temporal attention.

Generative Photography \citep{yuan2024genphoto} adds continuous control over camera parameters --
focal length, bokeh, shutter speed, colour temperature -- to a frozen text-to-video backbone, and
its architecture is reused in Chapter~\ref{cap:method}. A per-frame parameter value is turned into
a six-channel spatial conditioning tensor: three \emph{physical} channels that render the
parameter's effect, and three channels of a text-derived signal encoding the \emph{change} in the
parameter between consecutive frames (the construction its authors term \gls{ccl}). A small
encoder maps this tensor to features that are injected into the backbone's temporal attention
through zero-initialised \emph{merge} weights. This work substitutes an expression-intensity axis
for the camera axis, leaving the six-channel layout, the encoder, and the merge mechanism
unchanged; the substitution is developed in Chapter~\ref{cap:method}.

% =====================================================================
\section{Facial Expression Representation}\label{sec:bg-facs}

% \gls{facs} \citep{ekman1978facs}: anatomically grounded, per-AU, and
% crucially INTENSITY-GRADED, which is what makes a continuous control axis
% meaningful. Justify the choice of AU12 (lip-corner puller) as the single
% axis studied here: unambiguous, high dynamic range, reliably detectable.
% Mention the pre-diffusion lineage of AU-conditioned synthesis
% \citep{pumarola2018ganimation, choi2018stargan}.

The \gls{facs} \citep{ekman1978facs} decomposes any facial expression into \glspl{au}, the
movements of individual muscles or muscle groups, each scored not only as present or absent but on
a graded \emph{intensity} scale. This intensity coding is what makes a continuous control axis
meaningful: an expression is a point on a set of graded axes rather than a discrete label. This
work studies a single such axis, AU12 (the lip-corner puller), the defining action of a smile. AU12
is chosen because it is unambiguous, has a large and visible dynamic range, and is among the more
reliably detected action units, so it is a clean testbed for a continuous axis while the method
itself is not specific to it. Conditioning synthesis on action units predates diffusion: GAN-based
systems such as GANimation \citep{pumarola2018ganimation} and StarGAN \citep{choi2018stargan}
animate or translate expressions from action-unit or categorical labels, but without the
photorealism and text controllability of a diffusion backbone.

% =====================================================================
\section{Automatic Action Unit Estimation}\label{sec:bg-au-estimation}

% These tools are both the source of training labels and the measuring
% instrument, so their properties are part of the method's validity.
%   \citep{cheong2023pyfeat}      -- py-feat, used for labels AND evaluation
%   \citep{lugaresi2019mediapipe} -- MediaPipe blendshapes, independent check
%   \citep{chang2024libreface}    -- LibreFace, used by the prior art
%
% IMPORTANT: introduce the circularity problem HERE (same estimator produces
% labels and scores results), so that Cap04's mitigation reads as a designed
% choice rather than a defensive afterthought. All four related papers share
% this circularity without flagging it; this work flags and mitigates it.

Because no dataset provides per-frame ground-truth action-unit intensities at scale, the
intensities used here are estimated automatically. py-feat \citep{cheong2023pyfeat} predicts a
per-frame AU intensity and is used in this work both to \emph{label} the training data and to
\emph{score} the generated outputs. That double use is a validity concern: when the same estimator
produces the supervision and the metric, a model can score well simply by reproducing its
labeller's biases. The concern is raised here so that the mitigation in Chapter~\ref{cap:method}
reads as a designed choice rather than an afterthought -- a second estimator from a different
family, the face-mesh geometry of MediaPipe \citep{lugaresi2019mediapipe}, which produced none of
the labels, provides an independent check, and agreement between the two is what makes a result
trustworthy. Estimators of this kind, including LibreFace \citep{chang2024libreface}, are also used
by the prior art. All four systems of Chapter~\ref{cap:related-work} label and score with the same
estimator family without noting the circularity; flagging and mitigating it is part of this work's
evaluation.

% =====================================================================
\section{Datasets for Expression Intensity}\label{sec:bg-datasets}

% \gls{mead} \citep{wang2020mead}: multi-view, emotion-labelled, three
% discrete intensity levels, lab conditions. Explain what is extracted here
% that the discrete levels do not give: continuous per-frame AU12 from the
% onset-to-apex ramp of each clip.
% Contrast with the still-image AU datasets used by the prior art:
% DISFA \citep{mavadati2013disfa}, AffectNet \citep{mollahosseini2019affectnet}.
% Note the known consequence of lab-only training data (domain look,
% demographic bias) -- it comes back as a limitation in Cap05.

\gls{mead} \citep{wang2020mead} is a multi-view, emotion-labelled corpus of talking faces recorded
under controlled lighting, with three discrete emotion-intensity levels. The three levels are too
coarse for a continuous axis; what this work extracts instead is a \emph{continuous} per-frame AU12
value along the onset-to-apex ramp of each clip, so that a single clip supplies an ordered sequence
of smoothly increasing intensities. This differs from the still-image action-unit datasets used by
the prior art -- DISFA \citep{mavadati2013disfa}, with manually coded AU intensities, and AffectNet
\citep{mollahosseini2019affectnet}, a large in-the-wild collection -- which provide isolated frames
rather than intensity trajectories. Training on a single lab-recorded corpus carries a known cost
that returns as a limitation in Chapter~\ref{cap:preliminary-results}: the generated faces take on
the dataset's controlled-lighting, frontal appearance, and the model inherits its demographic skew.

% =====================================================================
\section{Evaluation Metrics}\label{sec:bg-metrics}

% Define each metric once, here, then only reference it in Cap04/Cap05.
%   - Pearson r between commanded and detected intensity (ramp following).
%   - \gls{cls}, the control-linearity protocol of \citet{hua2026pixelsmile}:
%     uniformly spaced commanded intensities, correlation against detected.
%   - Face-embedding cosine for identity \citep{deng2019arcface,
%     schroff2015facenet, cao2018vggface2}.
%   - \gls{lpips} \citep{zhang2018lpips} and CLIP-based prompt fidelity.
%   - Note that Pearson r is invariant to scale and offset, which is exactly
%     why an absolute-error metric is also required. This distinction is what
%     separates "follows the ramp" from "is calibrated", and it is the
%     central measurement point of Cap05.

Control is evaluated with a small set of metrics, defined here and only referenced later. Ramp
following is measured by the Pearson correlation $r$ between the commanded per-frame intensity and
the intensity detected in the generated frame. Absolute calibration is measured by the control
linearity score (\gls{cls}) of \citet{hua2026pixelsmile}: uniformly spaced commanded intensities
are correlated against the detected ones. Identity is measured by the cosine similarity of face
embeddings from a recognition network \citep{schroff2015facenet, cao2018vggface2, deng2019arcface};
image quality and prompt fidelity, where reported, use \gls{lpips} \citep{zhang2018lpips} and CLIP
similarity. One property of the correlation is central to Chapter~\ref{cap:preliminary-results}:
Pearson $r$ is invariant to scale and offset, so it rewards a model that \emph{orders and spaces}
the frames correctly even when the detected values are shifted or compressed. Distinguishing
``follows the ramp'' from ``is calibrated'' therefore requires an absolute-error metric -- the
\gls{mse} of the detected value against the commanded one -- alongside the correlation.
